\subsubsection{网络流经典模型}
\textbf{上下界无费用流}
\begin{enumerate}
	\item 无源汇上下界可行流：钦定容量下界流满并计算这种情况下每个点的净流入量 $in_u$，新建超级源点 $S$ 和超级汇点 $T$ 用以模拟净流入的影响。如果 $in_u>0$ 则连边 $(S,u,in_u)$，否则连边 $(u,T,-in_u)$。原图中其他边容量改为上下界之差。跑新图最大流，若满流则有可行流，把每条边原来的下界加上即可；否则无解。
	\item 有源汇上下界可行流：新建一条 $(T,S,+\infty)$ 的边即转化为无源汇上下界可行流。
	\item 有源汇上下界最大流：跑完无源汇的可行流后，断掉 $(T,S,+\infty)$ 的边后，残量网络上源汇之间（注意是原来的源汇不是新建的）继续跑最大流，加上即可。
	\item 有源汇上下界最小流：跑完无源汇的可行流后，断掉 $(T,S,+\infty)$ 的边后，残量网络上汇点到源点跑最大流，即求最大退流，减去即可。
\end{enumerate}
\textbf{上下界费用流}

以最小费用为例。若要求费用最大，全体费用取相反数以后做即可。鉴于可能出现负环，对于负权费用流后文会提到处理方式。以下默认费用为正。
\begin{enumerate}
	\item 无源汇上下界最小费用可行流：同无费用版本，把每一步都算上费用，最大流换成最小费用最大流。

	\textbf{注}：费用全正的情况下，不需要跑最小费用退流。如果能退，则存在权为负的增广路，没理由不在最小费用最大流的时候被处理过。
\item 有源汇上下界最小费用可行流：新建一条 $(T,S,+\infty,0)$ 的边即转化为无源汇上下界最小费用可行流。
\item 有源汇上下界最小费用最大流：跑完无源汇的最小费用可行流后，断掉 $(T,S,+\infty,0)$ 的边后，残量网络上源汇之间继续跑最小费用最大流，加上即可。
\item 有源汇上下界最小费用最小流：跑完无源汇的最小费用可行流后，断掉 $(T,S,+\infty,0)$ 的边后，残量网络上汇点到源点跑最小费用最大流，减去即可（费用是加上）。
\end{enumerate}
\textbf{负权费用流}

其实没那么神秘。对于每条负权边 $(u,v,c,w)$，我们把它变成两条边：一条是 $(u,v,c,w)$ 强制满流，另一条是 $(v,u,c,-w)$ 的退流边。强制满流直接当作上下界限制即可，在跑新图的时候负权边的容量会变成 0，等于没有这条边。这样所有费用都是正的。

\textbf{同余网络流}

现编的名字，即每条边流量都有一个全局的 $\bmod k$ 限制。把每条边的流量写成 $kx+y$，那么每个点处就会有关于 $x$ 的流量平衡，由于 $y$ 的净流入可能是 $k$ 的正数倍 $\Delta k$，所有流量平衡有一个偏差项 $\Delta$，与上下界一样用额外源汇模拟偏差即可在 $x$ 上做流子。

\textbf{最大权闭合子图}

建立源汇 $S,T$，对正权点连边 $(S,u,w_u)$，负权点连边 $(u,T,-w_u)$，原图边 $(u,v)$ 建 $(u,v,\infty)$，全图正权和减去最小割即为答案。

答案构造：在最大流结束的残量网络中，从源点沿正残量边可达的原图顶点组成所选闭合子图。

\textbf{最长反链}

对原图每个点拆点 $(in_u,out_u)$，原图边 $(u,v)$ 改为 $(out_u,in_v)$，所有 $in_u$ 点权为 -1，所有 $out_u$ 点权为 1。此图最大权闭合子图即为最长反链。

这个图比较好理解，因为一个最大的闭合子图一定是由若干个 $out_u$ 作为起点发出的（否则权值和非正），这几个 $u$ 构成了一个反链。

\paragraph{多测与注意} 每组清网络、需求平衡和基准费用。可行流后删超级源汇关联边及额外 $T\to S$ 边，正反向都删，先记录基准流。最小非负流的退流不超过基准流。最长反链需 DAG 的传递可达约束，可补 in 到 out 的无穷边；费用版记得加下界费用。
